\chapter{Wstęp}

Aplikacja \verb|Wydatkomat| została stworzona w celu uproszczenia zarządzania wspólnymi wydatkami pomiędzy znajomymi. System umożliwia kontrolowanie kosztów, tworzenie grup użytkowników, dodawanie wydatków oraz generowanie rozliczeń pomiędzy członkami grupy. Dzięki temu użytkownicy mogą w prosty sposób monitorować swoje zobowiązania finansowe oraz historię wydatków.

Projekt został zrealizowany jako aplikacja wieloplatformowa, dostępna z różnych urządzeń i systemów operacyjnych. System składa się z centralnego backendu udostępniającego REST API oraz klientów korzystających z tego API. Przygotowano:
\begin{itemize}
	\item aplikację webową stworzoną przy użyciu technologii React,
	\item aplikację mobilną opartą o React Native,
	\item aplikację desktopową napisaną w języku Python.
\end{itemize}

Backend aplikacji odpowiada za przetwarzanie danych, autoryzację użytkowników, logikę biznesową oraz komunikację z bazą danych. Komunikacja pomiędzy klientami a serwerem odbywa się przy użyciu protokołu HTTP oraz architektury REST.

Dodatkowym elementem systemu są statystyki wydatków, które pozwalają użytkownikom analizować swoje finanse oraz lepiej kontrolować budżet.

\section{Cel projektu}

Celem projektu było zaprojektowanie, implementacja, przetestowanie oraz wdrożenie wieloplatformowego systemu do zarządzania wspólnymi wydatkami. Projekt miał umożliwiać korzystanie z aplikacji na różnych urządzeniach przy zachowaniu wspólnego backendu oraz spójnej logiki działania systemu.

Istotnym założeniem projektu było również przygotowanie bezpiecznego i skalowalnego REST API, które może być wykorzystywane przez różne klienty aplikacji.

Aplikacja mobilna oraz desktopowa powinna umożliwiać również pracę w trybie offline, z późniejszą synchronizacją danych po ponownym połączeniu do sieci.

\section{Zakres projektu}

Zakres projektu obejmuje implementację systemu umożliwiającego:
\begin{itemize}
	\item rejestrację i uwierzytelnianie użytkowników,
	\item zarządzanie swoim kontem użytkownika,
	\item zarządzanie znajomymi i uczestnikami wydatków,
	\item dodawanie i edytowanie wydatków,
	\item podział kosztów pomiędzy użytkownikami,
	\item przeglądanie historii wydatków,
	\item generowanie statystyk finansowych,
	\item zarządzanie użytkownikami jako adminsitrator,
	\item otrzymywanie powiadomień wybranymi przez siebie kanałami.
\end{itemize}

Projekt obejmuje również przygotowanie dokumentacji REST API zgodnej ze standardem OpenAPI oraz wdrożenie aplikacji w środowisku zdalnym.